% Preamble of "Numerical Derivatives with Numerixx" -- the layout and design language of the EquiCore
% design paper (Tufte layout after the Typst template toffee-tufte): a wide right margin for sidenotes and
% margin floats, full-width blocks, IBM Plex Sans with sans-serif math, bold numbered headings, block
% paragraphs, dark-gray text, one blue accent, and code on a light-gray ground. Changes from that paper,
% all for a book: the report class with chapters and parts, a numbered sidebar box in the same design
% language, a list of sidebars, and chapter-level running heads.
%
% LAYOUT VOCABULARY FOR AUTHORS (the same as the EquiCore paper, plus sidebars)
%   \sidenote{...}          numbered note in the right margin (use instead of \footnote)
%   \mnote{...}             unnumbered margin note (a short gloss or definition); NEVER \marginnote in running
%                           text (it does not avoid collisions; the marginnote package is kept for \sidecaption)
%   figure / table          text-width float (text block 114 mm); caption in the margin with \sidecaption{...}
%                           (label inside: \sidecaption{...\label{...}}), issued at the TOP of the float
%   figure* / table*        full-width float over text block and margin (174 mm = \fullwidthlen); ordinary \caption
%   marginfigure / margintable   small float in the margin (53 mm); use \caption{...}
%   \begin{fullwidth} ... \end{fullwidth}   full-width running material (a long listing, a wide equation)
%   \begin{sidebar}[label=sb:...]{Title} ... \end{sidebar}   (the label as an option, not a \label inside)
%                           a numbered box (Sidebar 3.2) for a complex topic a user does not need to follow;
%                           the text around it must read without it. Text width; it may break across pages.
%   Citations: numbered, bracketed superscripts (\cite); reference list at the end (fullbib).
%   Plots: \begin{axis}[eqtufte, ...] in every pgfplots figure, data from data/*.csv only.
%   Code: \lstinputlisting[linerange=...]{code/<file>.cpp} from compiled examples only.
%
% ---------------------------------------------------------------- page
\usepackage[a4paper,
  left=18mm, top=25mm, bottom=25mm,
  textwidth=114mm, marginparsep=7mm, marginparwidth=53mm,
  headheight=20pt, headsep=8mm, footskip=11mm]{geometry}
\newlength{\fullwidthlen}\setlength{\fullwidthlen}{\dimexpr\textwidth+\marginparsep+\marginparwidth\relax}
% ---------------------------------------------------------------- fonts
\usepackage[T1]{fontenc}
\usepackage[sfdefault]{plex-sans}                  % IBM Plex Sans, the text face of toffee-tufte
\usepackage[scale=0.92]{plex-mono}                 % IBM Plex Mono for code
\usepackage{amsmath}
\usepackage[amssymbols,slantedGreek]{newtxsf}      % sans-serif math (STIX sans) to match the text
\usepackage{microtype}
\DisableLigatures[<,>]{encoding = T1, family = tt*}  % no guillemets from >> in code
\usepackage{setspace}
\setstretch{1.12}
% ---------------------------------------------------------------- core packages
\usepackage{bm}
\usepackage{booktabs,multirow,array,tabularx,threeparttable}
\usepackage{multicol}
\usepackage{flafter}    % a float never appears before its place in the text
\usepackage{graphicx}
\usepackage{xcolor}
\definecolor{accent}{HTML}{0B5FAE}   % the one accent: links, cross-references and highlighted data
\definecolor{muted}{gray}{0.42}
\definecolor{ink}{gray}{0.15}        % text: dark gray, not black (toffee-tufte: luma 15 %)
\definecolor{codebg}{gray}{0.95}     % ground of code blocks, inline code and sidebars
\usepackage[per-mode=symbol,separate-uncertainty=true,detect-all]{siunitx}
\usepackage{tikz}
\usetikzlibrary{arrows.meta,positioning,shapes.geometric,fit,calc,backgrounds}
\usepackage{pgfplots}\pgfplotsset{compat=1.18}
\usepgfplotslibrary{groupplots}
% Tufte-style plots: \begin{axis}[eqtufte, ...] in every pgfplots figure (thin gray axes, sans-serif labels,
% no frame, no grid; label curves directly instead of a legend where possible).
\pgfplotsset{eqtufte/.style={axis lines=left, axis line style={thin,black!55,-},
  tick style={thin,black!55}, major tick length=2pt, minor tick length=1pt,
  tick label style={font=\sffamily\scriptsize,black!70}, label style={font=\sffamily\scriptsize},
  title style={font=\sffamily\scriptsize}, legend style={draw=none,fill=none,font=\sffamily\scriptsize},
  every axis plot/.append style={line width=0.6pt}, mark size=0.7pt, clip=false,
  enlargelimits=false, scale only axis}}
\usepackage[ruled,vlined,linesnumbered]{algorithm2e}
\usepackage{listings}
% ---------------------------------------------------------------- paragraphs
\usepackage[skip=0.55\baselineskip plus 2pt minus 1pt, indent=0pt]{parskip}
% ---------------------------------------------------------------- captions, margin material
\usepackage[font={small},labelfont={small,bf},labelsep=period,
  justification=raggedright,singlelinecheck=false,skip=4pt]{caption}
\usepackage[oneside]{sidenotes}
\usepackage{marginfix}
\renewcommand*{\marginfont}{\footnotesize\normalfont\setstretch{1.0}}
\DeclareCaptionStyle{marginfigure}{font={footnotesize},labelfont={footnotesize,bf},justification=raggedright}
\DeclareCaptionStyle{margintable}{font={footnotesize},labelfont={footnotesize,bf},justification=raggedright}
\DeclareCaptionStyle{sidecaption}{font={footnotesize},labelfont={footnotesize,bf},justification=raggedright}
\DeclareCaptionStyle{widefigure}{font={small},labelfont={small,bf},justification=raggedright}
\DeclareCaptionStyle{widetable}{font={small},labelfont={small,bf},justification=raggedright}
\newcommand{\mnote}[1]{\marginpar{\marginfont\raggedright #1}}
\makeatletter
\RenewDocumentCommand\@sidenotes@placemarginal{+m +m}{%
  \IfNoValueOrEmptyTF{#1}{\marginpar{\marginfont\raggedright #2}}{\marginnote{\marginfont\raggedright #2}[#1]}}
\makeatother
\newenvironment{fullwidth}{\begin{adjustwidth}{}{-\dimexpr\marginparsep+\marginparwidth\relax}}{\end{adjustwidth}}
% ---------------------------------------------------------------- headings (bold, numbered)
\usepackage{titlesec}
\setcounter{secnumdepth}{2}
\setcounter{tocdepth}{1}
% Parts: a page of their own, the part number muted above the title.
\titleformat{\part}[display]{\normalfont\filright\bfseries}
  {\fontsize{11}{13}\selectfont\color{muted}Part \thepart}{0.8ex}{\fontsize{22}{27}\selectfont}
% Chapters: a new page, the chapter number muted above a large bold title.
\titleformat{\chapter}[display]{\normalfont\bfseries}
  {\fontsize{11}{13}\selectfont\color{muted}\csname @chapapp\endcsname\ \thechapter}{0.6ex}{\fontsize{18}{22}\selectfont}
\titlespacing*{\chapter}{0pt}{0pt}{4ex}
\titleformat{\section}{\normalfont\fontsize{13}{16}\selectfont\bfseries}{\thesection}{0.6em}{}
\titleformat{\subsection}{\normalfont\fontsize{11}{14}\selectfont\bfseries}{\thesubsection}{0.6em}{}
\titleformat{\paragraph}[runin]{\normalfont\bfseries}{}{0pt}{}  % write the period in the heading text
\titlespacing*{\section}{0pt}{2.4ex plus .6ex minus .2ex}{1.0ex}
\titlespacing*{\subsection}{0pt}{2.0ex plus .5ex minus .2ex}{0.7ex}
\titlespacing*{\paragraph}{0pt}{0.9ex plus .3ex}{0.6em}
% ---------------------------------------------------------------- running heads
\usepackage{fancyhdr}
\usepackage{lastpage}
\newcommand{\bkRunningTitle}{Numerical derivatives with Numerixx}
\newcommand{\bkRunningAuthor}{Kenneth Troldal Balslev}
\newcommand{\bkRunningDate}{Draft}
\pagestyle{fancy}\fancyhf{}
\renewcommand{\headrulewidth}{0pt}
\renewcommand{\chaptermark}[1]{\markboth{\thechapter.\ #1}{}}
\fancyhead[L]{\footnotesize\itshape\color{muted}\bkRunningTitle\\\leftmark}
\fancyhead[R]{\footnotesize\itshape\color{muted}\bkRunningAuthor\\\bkRunningDate}
\fancyfoot[R]{\footnotesize\itshape\color{muted}\thepage/\pageref*{LastPage}}
\fancypagestyle{plain}{\fancyhf{}\fancyfoot[R]{\footnotesize\itshape\color{muted}\thepage/\pageref*{LastPage}}}
\fancyheadoffset[R]{\dimexpr\marginparsep+\marginparwidth\relax}
\fancyfootoffset[R]{\dimexpr\marginparsep+\marginparwidth\relax}
% ---------------------------------------------------------------- title page
% Title, subtitle, author and edition centered over the full width (text block + margin).
\makeatletter
\newcommand{\bkSubtitle}{}
\renewcommand{\maketitle}{%
  \begin{titlepage}\color{ink}\vspace*{0.22\textheight}
  \begin{fullwidth}\centering
    {\fontsize{24}{29}\selectfont\bfseries\@title\par}\vspace{2ex}
    {\fontsize{13}{16}\selectfont\color{muted}\bkSubtitle\par}\vspace{6ex}
    {\large\@author\par}\vspace{1ex}
    {\normalsize\color{muted}\@date\par}
  \end{fullwidth}\end{titlepage}}
\makeatother
% ---------------------------------------------------------------- sidebars
% A numbered box, numbered within the chapter (Sidebar 3.2), on the code ground with a thin accent rule on the
% left, the label and title in bold. Text width; breakable; listed in the List of sidebars. The optional argument
% takes tcolorbox options, above all label=sb:... (a \label inside the box would refer to the enclosing section).
\usepackage[most]{tcolorbox}
\newtcolorbox[auto counter, number within=chapter, list inside=sbr,
  list type=section]{sidebar}[2][]{%
  enhanced, breakable, colback=codebg, colframe=codebg, boxrule=0pt, arc=0pt, outer arc=0pt,
  borderline west={1.2pt}{0pt}{accent}, left=7pt, right=7pt, top=5pt, bottom=5pt,
  before skip=1.1\baselineskip, after skip=1.1\baselineskip,
  fontupper=\small\setstretch{1.08}, parbox=false,
  title={Sidebar~\thetcbcounter.\ #2}, fonttitle=\small\bfseries, coltitle=ink, detach title,
  before upper={\tcbtitle\par\smallskip},
  list entry={\protect\numberline{\thetcbcounter}#2}, #1}
\newcommand{\listofsidebars}{\tcblistof[\chapter*]{sbr}{List of sidebars}}
% ---------------------------------------------------------------- references
\usepackage[journal=iecred,doi=true]{achemso}  % numbered reference list (package form; loads natbib + mciteplus)
\setcitestyle{super,comma,sort&compress,open={[},close={]}}  % bracketed superscripts, never confused with sidenotes
\usepackage[colorlinks=true,linkcolor=accent,citecolor=accent,urlcolor=accent,filecolor=accent]{hyperref}
\usepackage[capitalise,nameinlink]{cleveref}
\crefname{tcb@cnt@sidebar}{Sidebar}{Sidebars}
\Crefname{tcb@cnt@sidebar}{Sidebar}{Sidebars}
% ---------------------------------------------------------------- tables
\setlength{\heavyrulewidth}{0.08em}
\setlength{\lightrulewidth}{0.05em}
\setlength{\cmidrulewidth}{0.03em}
% ---------------------------------------------------------------- code listings (C++)
\lstdefinestyle{bkcode}{language=C++,basicstyle=\ttfamily\small,
  keywordstyle=\color{accent},commentstyle=\itshape\color{muted},
  stringstyle=\color{muted},columns=fullflexible,keepspaces=true,
  breaklines=true,frame=single,framerule=0pt,rulecolor=\color{codebg},
  backgroundcolor=\color{codebg},framesep=5pt,xleftmargin=5pt,xrightmargin=5pt,
  aboveskip=0.8\baselineskip,belowskip=0.6\baselineskip,
  morekeywords={concept,requires,constexpr,consteval,noexcept,auto,std}}
\lstset{style=bkcode}
% \snippet{<file>}{<name>}: the lines of code/<file>.cpp between the comment lines "// <name>" and "// </name>",
% without the markers. Markers and snippet start in column 0 (\lstinputlisting cannot gobble indentation, and an
% indented end marker is not recognised); every file in code/ follows the convention (code/README.md).
% The range ends are braced: without the braces, linerange splits a hyphenated name such as first-ridders at its
% first hyphen and quotes the wrong lines, or nothing, without a warning (tools/check.py checks the markers).
\newcommand{\snippet}[2]{\lstinputlisting[rangeprefix=//\ <, rangesuffix=>, linerange={{#2}-{/#2}},
  includerangemarker=false]{code/#1.cpp}}
% ---------------------------------------------------------------- float placement
\renewcommand{\topfraction}{0.85}
\renewcommand{\bottomfraction}{0.6}
\renewcommand{\textfraction}{0.12}
\renewcommand{\floatpagefraction}{0.7}
\setcounter{topnumber}{3}
\setcounter{totalnumber}{4}
\makeatletter
\setlength{\@fptop}{0pt}
\setlength{\@fpsep}{16pt plus 6pt}
\setlength{\@fpbot}{0pt plus 1fil}
\makeatother
% ---------------------------------------------------------------- line breaking, widows, clubs
\setlength{\emergencystretch}{2.5em}
\widowpenalty=10000
\clubpenalty=10000
% ---------------------------------------------------------------- text color
\AtBeginDocument{\color{ink}}
% ---------------------------------------------------------------- reference list
\usepackage{needspace}
\newenvironment{fullbib}%
  {\chapter*{\refname}\addcontentsline{toc}{chapter}{\refname}\begin{fullwidth}\footnotesize\setstretch{1.0}%
   \renewcommand{\bibsection}{}\setlength{\columnsep}{8mm}\raggedcolumns\raggedright\begin{multicols}{2}}%
  {\end{multicols}\end{fullwidth}}
